% GENERATED FILE. Edit canonical lessons, then run npm run generate:books.
% canonical-source-hash: fnv1a64:a7bf1fda823c9c5e
% canonical-lessons: PT-C403-vidro, PT-C403-retrovisor, PT-C403-buzina, PT-C403-buzinar, PT-C403-embaciar

\chapter{Vidros e espelhos}
\label{ch:pt-vidros-e-espelhos}
\hlchaptermodality{403}

\begin{chapteropening}
\textbf{By the end of this chapter:} \emph{I can talk about seeing clearly from a car: the glass, the mirror, the horn, honking and windows that mist up.}

Complete the last lesson of chapter 403: I can talk about seeing clearly from a car: the glass, the mirror, the horn, honking and windows that mist up.
\end{chapteropening}

\section[o vidro]{o vidro — the glass (material); the car window}

\label{lesson:PT-C403-vidro}



\begin{quote}
Before the new one: say the Portuguese for the clutch, then the Portuguese for to put into gear.
\end{quote}

\subsection*{You'll want to know: o vidro}
\textbf{o vidro} — \textquotedblleft{}the glass (material); the car window\textquotedblright{}.

\begin{quote}
Ela abre o vidro do carro porque está calor.
\end{quote}

\begin{quote}
\emph{She opens the car window because it is hot.}
\end{quote}

\emph{Uma garrafa de vidro} — a glass bottle. A drinking glass is \emph{um copo}.

\begin{cousinweb}
\emph{vidro} is inherited from Latin \emph{vitrum} 'glass' (EtymDB).
\end{cousinweb}

\begin{grammarlens}[title={What you've built}]
One more word for telling, explaining and describing at B1.
\end{grammarlens}

\subsection*{Your turn}
\begin{itemize}
\raggedright
  \item \emph{Say it:} \emph{o vidro}
  \item \emph{Say it:} \emph{o vidro}, once more
  \item \emph{Say it:} a sentence of your own with \emph{o vidro}
  \item \emph{Recall:} say the Portuguese for the clutch, then the Portuguese for to put into gear, then say \emph{o vidro} again
\end{itemize}

\begin{tcolorbox}[breakable,colback=teal!4,colframe=teal!35!black,title={Before you move on}]
What does o vidro mean? (\textquotedblleft{}The glass (material); the car window\textquotedblright{}.) Say it once more.
\end{tcolorbox}
\section[o retrovisor]{o retrovisor — the rear-view mirror, the wing mirror}

\label{lesson:PT-C403-retrovisor}



\begin{quote}
Before the new one: say the Portuguese for to put into gear, then the Portuguese for the glass.
\end{quote}

\subsection*{You'll want to know: o retrovisor}
\textbf{o retrovisor} — \textquotedblleft{}the rear-view mirror, the wing mirror\textquotedblright{}.

\begin{quote}
Antes de mudar de faixa, ele olha para o retrovisor.
\end{quote}

\begin{quote}
\emph{Before changing lanes, he checks the mirror.}
\end{quote}

\emph{O espelho retrovisor} in full.

\begin{cousinweb}
\emph{retrovisor} is a learned formation from \emph{retro-} 'backwards' and \emph{visor} 'viewer', built on Latin \emph{vidēre} 'to see' (Lewis \& Short): the mirror that sees behind.
\end{cousinweb}

\begin{grammarlens}[title={What you've built}]
One more word for telling, explaining and describing at B1.
\end{grammarlens}

\subsection*{Your turn}
\begin{itemize}
\raggedright
  \item \emph{Say it:} \emph{o retrovisor}
  \item \emph{Say it:} \emph{o retrovisor}, once more
  \item \emph{Say it:} a sentence of your own with \emph{o retrovisor}
  \item \emph{Recall:} say the Portuguese for to put into gear, then the Portuguese for the glass, then say \emph{o retrovisor} again
\end{itemize}

\begin{tcolorbox}[breakable,colback=teal!4,colframe=teal!35!black,title={Before you move on}]
What does o retrovisor mean? (\textquotedblleft{}The rear-view mirror, the wing mirror\textquotedblright{}.) Say it once more.
\end{tcolorbox}
\section[a buzina]{a buzina — the horn (of a car)}

\label{lesson:PT-C403-buzina}



\begin{quote}
Before the new one: say the Portuguese for the glass, then the Portuguese for the rear-view mirror.
\end{quote}

\subsection*{You'll want to know: a buzina}
\textbf{a buzina} — \textquotedblleft{}the horn (of a car)\textquotedblright{}.

\begin{quote}
O motorista do táxi tocou a buzina duas vezes.
\end{quote}

\begin{quote}
\emph{The taxi driver sounded the horn twice.}
\end{quote}

\emph{Tocar a buzina} — to sound the horn.

\begin{cousinweb}
\emph{buzina} is inherited from Latin \emph{būcina} 'curved trumpet, horn' (Lewis \& Short); EtymDB lists the Portuguese word without a source.
\end{cousinweb}

\begin{grammarlens}[title={What you've built}]
One more word for telling, explaining and describing at B1.
\end{grammarlens}

\subsection*{Your turn}
\begin{itemize}
\raggedright
  \item \emph{Say it:} \emph{a buzina}
  \item \emph{Say it:} \emph{a buzina}, once more
  \item \emph{Say it:} a sentence of your own with \emph{a buzina}
  \item \emph{Recall:} say the Portuguese for the glass, then the Portuguese for the rear-view mirror, then say \emph{a buzina} again
\end{itemize}

\begin{tcolorbox}[breakable,colback=teal!4,colframe=teal!35!black,title={Before you move on}]
What does a buzina mean? (\textquotedblleft{}The horn (of a car)\textquotedblright{}.) Say it once more.
\end{tcolorbox}
\section[buzinar]{buzinar — to honk, to sound the horn}

\label{lesson:PT-C403-buzinar}



\begin{quote}
Before the new one: say the Portuguese for the rear-view mirror, then the Portuguese for the horn.
\end{quote}

\subsection*{You'll want to know: buzinar}
\textbf{buzinar} — \textquotedblleft{}to honk, to sound the horn\textquotedblright{}.

\begin{quote}
Os carros buzinam quando o semáforo fica verde.
\end{quote}

\begin{quote}
\emph{The cars honk when the lights turn green.}
\end{quote}

Built on \emph{a buzina}, the horn.

\begin{cousinweb}
\emph{buzinar} is built on \emph{buzina} 'horn', which is inherited from Latin \emph{būcina} 'trumpet' (Lewis \& Short).
\end{cousinweb}

\begin{grammarlens}[title={What you've built}]
One more word for telling, explaining and describing at B1.
\end{grammarlens}

\subsection*{Your turn}
\begin{itemize}
\raggedright
  \item \emph{Say it:} \emph{buzinar}
  \item \emph{Say it:} \emph{buzinar}, once more
  \item \emph{Say it:} a sentence of your own with \emph{buzinar}
  \item \emph{Recall:} say the Portuguese for the rear-view mirror, then the Portuguese for the horn, then say \emph{buzinar} again
\end{itemize}

\begin{tcolorbox}[breakable,colback=teal!4,colframe=teal!35!black,title={Before you move on}]
What does buzinar mean? (\textquotedblleft{}To honk, to sound the horn\textquotedblright{}.) Say it once more.
\end{tcolorbox}
\section[embaciar]{embaciar — to mist up, to steam up}

\label{lesson:PT-C403-embaciar}



\begin{quote}
Before the new one: say the Portuguese for the horn, then the Portuguese for to honk.
\end{quote}

\subsection*{You'll want to know: embaciar}
\textbf{embaciar} — \textquotedblleft{}to mist up, to steam up\textquotedblright{}.

\begin{quote}
Com o frio, os vidros embaciam por dentro.
\end{quote}

\begin{quote}
\emph{In the cold, the windows mist up on the inside.}
\end{quote}

\emph{Os vidros estão embaciados} — the windows are misted up.

\begin{cousinweb}
\emph{embaciar} is built from \emph{em-} and \emph{baço} 'dull, without shine'. \emph{Baço} is inherited from Latin \emph{badius} 'reddish-brown' (EtymDB).
\end{cousinweb}

\begin{grammarlens}[title={What you've built}]
One more word for telling, explaining and describing at B1.
\end{grammarlens}

\subsection*{Your turn}
\begin{itemize}
\raggedright
  \item \emph{Say it:} \emph{embaciar}
  \item \emph{Say it:} \emph{embaciar}, once more
  \item \emph{Say it:} a sentence of your own with \emph{embaciar}
  \item \emph{Recall:} say the Portuguese for the horn, then the Portuguese for to honk, then say \emph{embaciar} again
\end{itemize}

\begin{tcolorbox}[breakable,colback=teal!4,colframe=teal!35!black,title={Before you move on}]
What does embaciar mean? (\textquotedblleft{}To mist up, to steam up\textquotedblright{}.) Say it once more.
\end{tcolorbox}
